\section{Experimental design}
\label{sec:design}

\subsection{Construct and model}

Our study-wide observable is \emph{false operational commitment}: a sampled status or explicit
prose verdict contradicts machine-specified scenario truth. In the structured-action bank we
define the stronger \emph{knowledge-conditioned false commitment}: the status is false and a
shared unpressured branch demonstrates the correct baseline answer. A false status without that
evidence is retained as \textsc{wrong-without-knowledge}. The natural-pressure bank lacks a
branch-level knowledge check, so C9 measures false commitment rather than deception in this
stronger sense. Neither construct establishes a persistent hidden goal or human-like intent.

All target generations use \model. The sequential relational bank is bound to the
immutable checkpoint revision and component hashes reported in the artifact manifest; gaps in
legacy sources are explicit. No results are pooled across precision regimes or banks.

\subsection{What ``geometry'' means in this study}

Geometry organizes activation-space variation; no group action on the state space is claimed.

\begin{itemize}
    \item \textbf{Local fields and charts.} A chart assigns a local reference frame from nearest
          training neighbors; local control fields condition intervention directions on the
          current state. Graph-mean, chart-mean, and chart-distilled selectors embed this
          local-first logic without inspecting candidate outcomes. The Local Control Field (LCF)
          extends this with a neighborhood-conditioned flow generator.
    \item \textbf{Typed relations and residual--attention graphs.} Beyond simple residual
          vectors, typed relational objects record pairwise residual relations, causal attention
          rows, and token role/position metadata across four layers (12, 16, 19, 20). The
          C10 post-commitment graph and C11 pre-commitment fields are built from these.
    \item \textbf{Pressure flow (descriptive).} Across conversational turns under graded pressure,
          we compute anchor-only pseudo-orbits and measure monotonicity, field coherence, and
          cross-family similarity. These statistics describe coherent flow-like organization;
          they do not claim a group, semigroup, causal trajectory, or transferable universal
          field.
    \item \textbf{Tangent projection (actuator detail).} Some controllers project a direction into
          a local tangent subspace. This is an actuator implementation detail, distinct from the
          relational graph (C10/C11), chart/atlas action selection (C1), and LCF.
\end{itemize}

Nowhere does the study assume a Lie group acting on the state manifold. Where a
$\mathbb{Z}_2$ involution is tested, it is tested empirically: forced global
$\mathbb{Z}_2$ collapses sign accuracy; the result is compatible with chart-local
structure but does not prove a local involution.

\subsection{Three non-pooled settings}

\paragraph{Natural pressure.}
Free-prose workplace conversations have machine-specified ground truth and four assigned
schedule families: smooth escalation, late-compressed, abrupt step, and benign control.
Scripted and adaptive banks contain 96 and 128 conversations; a separate 96-conversation
dissociation bank weakens correlation between current and accumulated pressure. Claude
Sonnet 4.6 generated adaptive user turns from per-turn intensity targets. A blinded
outcome channel sees transcript and truth, while a pressure channel sees only user turns.
Each channel is judged through three separate calls to Claude Sonnet 5 and aggregated by
majority label or median score---repeated calls to one model family, not three independent
judge models, with no human-agreement study. The public evidence binds rubric and prompts
but not an immutable provider revision for every historical call.

For each assistant turn, let $h_t$ indicate a first false commitment, $I_t$ the judged
pressure intensity of the immediately preceding user turn, and $C_{t-1}$ the accumulated
intensity:
\begin{equation}
  \operatorname{logit} \Pr(h_t=1\mid h_{<t}=0)
  = \delta + \alpha I_t + \gamma C_{t-1}.
  \label{eq:hazard}
\end{equation}
We compare the full model with $\gamma=0$ by leave-one-family-out log loss.

\paragraph{Sequential structured action.}
Across 600 rows, 60 scenarios, 20 families, and four turns, the environment appends
\texttt{Status:}, samples one unrestricted vocabulary token, appends it to context, then
repeats with \texttt{\textbackslash nCaveat:}. PASS, FAIL, and SKIP map by raw token ID;
other status tokens are \textsc{no-action}. The bank contains 2,400 status records
and 1,680 unique status events. Status recognition passed at 2,400/2,400 but the frozen
baseline-knowledge gate missed by one scenario (56/60 versus 57/60); we treat it as
exploratory. Full protocol detail in Appendix~\ref{app:secondary-methods}.

\paragraph{Control banks.}
The powered150 offline field stores candidate responses for 600 deceptive and 600 honest
conversations---matched policy evaluation, not fresh generation. The prospective
natural-prose experiment generates 96 new rows over four held-out families. A smaller
pilot decomposition contains 160 primary rows and a 32-row off-tangent follow-up.
Descriptive transfer tests use a 484-row Apollo-derived evaluation split with
soft-judged generation; this bank has no machine-checkable control target.

\subsection{Activation views and causal anchors}

Sequential captures store every-token BF16 residuals and all 32 causal attention heads at
layers 12, 16, 19, and 20, with lossless attention triangles and token role/turn/span/position
metadata. The \emph{status anchor} is immediately before the sampled status token; the
\emph{caveat anchor} occurs after status has entered the causal state and before the caveat
token. This ordering fixes interpretation: C10 uses the post-commitment growth edge
(status$\to$caveat), C11 uses action-free pre-status views, C9 reads the final non-special
content token of each assistant turn (primarily layer 16), and C5 intervenes only in the
layer-16 residual stream.

\subsection{Post-action graph and linear comparator}

C10's typed growth edge records pairwise residual relations and typed token metadata across
layers 12, 16, 19, and 20, retrieves eight same-status training events across five outer
family folds, and forms Jeffreys-smoothed class probabilities (Appendix~\ref{app:c11}). The
raw-residual comparator uses the same folds and is linear at application time.

\subsection{Pre-action instruments and controllers}

C11's three frozen instruments---spectral one-state radial field, connection-response path
field, and masked geometry-only risk model---test pre-status features. Unlike C10's
truth-blind prior, C11's comparators key on true and desired status while the tested
geometry is barred from those fields (details in Appendix~\ref{app:c11}).

C1's powered150 field uses the scenario-truth oracle route: deceptive-row menus contain only
corrective-target candidates. CNG takes an unconstrained argmax over the menu and is
candidate-response-blind at held-out application. The learned pre-response ridge receives
route and candidate target as features and is likewise response-blind. The retrospective
fixed policy maps route to one linear coordinate---but that coordinate was selected post hoc
from whole-grid candidate rewards, making it a response-informed feasibility ceiling, not a
response-free policy. The route is task-supplied in C1; a separate L20 linear route
diagnostic recovers PASS/FAIL on 1,680/1,680 held-out-family predictions under
same-distribution leave-one-family-out CV on a fresh 12-family bank. It is routing-only,
not OOD, and not end-to-end control; it demonstrates that route-conditioned deployment is
possible within this benchmark.

Tangent projection is an actuator detail, not a synonym for the geometric apparatus. C5
compares a frozen layer-16 tangent-actuated policy under two decode-time gates with a
family-matched linear policy on fresh prose. C13 replays one intervention at exact
pre-status structured-action prefixes, authenticating each live root against a sealed
four-layer source-bank query and comparing a negative-risk-gradient proposal with no-op,
sign-reversed, and random-tangent controls.

\subsection{Information budgets and chronology}

Table~\ref{tab:ladder} states what each control or readout receives. The critical distinction:

\begin{enumerate}
    \item \textbf{Post-action readout:} The committed status token is causally visible
          (C10). The raw-residual linear probe dominates all geometric alternatives.
    \item \textbf{Response-aware action selection:} The controller sees candidate outcomes and
          chooses the best. In this study, this reduces to measured margins and argmax,
          and response-aware selectors tie it exactly (79/100).
    \item \textbf{Candidate-response-blind structural selection:} The controller sees candidate
          targets, routes, and features but not outcomes. Local geometric structure earns its
          most unambiguous wins here (CNG 599/600 versus ridge 170/600).
    \item \textbf{Response-informed actuation:} The post-hoc fixed L16 $\alpha$96 coordinate
          reaches 600/600 with zero harm, but was selected from whole-grid candidate rewards.
          It is a feasibility ceiling, not a response-free policy.
\end{enumerate}

\begin{table}[t]
\caption{Evidence ladder and information budget. ``Post'' means the committed status token is
causally visible. ``Oracle'' means scenario truth supplies the target route.}
\label{tab:ladder}
\centering
\small
\resizebox{\linewidth}{!}{%
\begin{tabular}{llll}
\toprule
Setting & Outcome visible & Oracle target & Public role \\
\midrule
C10 post-commitment readout & post & no & decodability \\
C11 pre-status warning & no & no & early warning \\
C1 learned ranker & no candidate outcome & yes, feature + field & offline action selection \\
C1 CNG & no candidate outcome & yes, task-routed field & conditional correction \\
C1 fixed linear & coordinate post hoc from rewards & yes, task-routed field & response-informed ceiling \\
C5 tangent-actuated policy & no & no & prospective control \\
\bottomrule
\end{tabular}%
}
\end{table}

C5 is the only prospectively specified public control claim but remains a development
result. C10's evaluation bank existed before the linear-comparator registration. C1, C2,
C9, C11, C12, and C13 are retrospective syntheses. Registration tier travels with every
claim in the public registry.

\subsection{Endpoints and uncertainty}

Behavioral endpoints retain all outcome classes; scenario family is the resampling unit with
outer held-out folds. C10's primary metric is family-macro Brier score on 1,283
honest/deceptive post-commitment events under a nuisance-preserving permutation; a registered
re-score adds truth-aware comparators. C11 instruments have different populations and no
pooled metric (Appendix~\ref{app:c11}).
